\documentclass[12pt,a4paper]{article}
\usepackage[portuguese]{babel}
\usepackage[utf8]{inputenc}
\usepackage[T1]{fontenc}
\usepackage{graphicx}
\usepackage{hyperref}
\usepackage{geometry}
\usepackage{fancyhdr}
\usepackage{titlesec}
\usepackage{listings}
\usepackage{xcolor}

\geometry{margin=2.5cm}

\hypersetup{
    colorlinks=true,
    linkcolor=blue,
    filecolor=magenta,
    urlcolor=blue,
}

\lstset{
    backgroundcolor=\color{lightgray!20},
    basicstyle=\ttfamily\small,
    breaklines=true,
    frame=single,
    captionpos=b,
    tabsize=2
}

\pagestyle{fancy}
\fancyhf{}
\rhead{\thepage}
\lhead{Sistema de Alarme Doméstico}
\cfoot{Copyright © 2026 - Ana Raquel}

\titleformat{\section}{\Large\bfseries}{\thesection}{1em}{}
\titleformat{\subsection}{\large\bfseries}{\thesubsection}{1em}{}

\begin{document}

\begin{titlepage}
    \centering
    \vspace*{2cm}

    {\LARGE\bfseries Sistema de Alarme Doméstico\par}
    \vspace{1cm}
    {\Large Manual de Utilização\par}
    \vspace{2cm}

    {\large\bfseries Projeto de Programação Orientada aos Objetos\par}
    \vspace{3cm}

    {\Large Ana Raquel\par}
    \vspace{1cm}
    {\large 2026\par}

    \vspace{3cm}
    \includegraphics[width=0.3\textwidth]{logo.png}

    \vfill
\end{titlepage}

\newpage
\tableofcontents
\newpage

\section{Introdução}

Este documento descreve o sistema de alarme doméstico desenvolvido no âmbito da
uc de Programação Orientada aos Objetos. O sistema simula um alarme para um
apartamento com cinco divisões, gerindo sensores de diferentes tipos e
controlando o estado do alarme através de uma senha.

\subsection{Objetivos do Sistema}

O sistema de alarme foi desenvolvido para cumprir os seguintes requisitos:

\begin{itemize}
    \item Gerir o estado do alarme (armado/desarmado)
    \item Autenticar utilizadores através de senha de 4 dígitos
    \item Controlar uma campainha de alarme
    \item Processar eventos de sensores de intrusão, movimento, fogo, gás e inundação
    \item Executar simulações baseadas em ficheiros de eventos
    \item Oferecer uma interface de menu interativo
\end{itemize}

\subsection{Visão Geral}

O sistema é composto por três classes principais:

\begin{itemize}
    \item \textbf{ComandoAlarme}: Gere o estado do alarme, autenticação e campainha
    \item \textbf{Sensor}: Representa os sensores nas divisões do apartamento
    \item \textbf{Simulacao}: Coordena a simulação e processamento de eventos
\end{itemize}

\newpage

\section{Arquitetura do Sistema}

\subsection{Diagrama de Classes}

O sistema é composto pelas seguintes classes principais:

\subsubsection{Classe ComandoAlarme}

\begin{figure}[h!]
    \centering
    \begin{tabular}{|l|p{10cm}|}\hline
        \textbf{Atributos} & \textbf{Descrição}\\\hline
        estado & Boolean que indica se o sistema está armado ou desarmado\\\hline
        senha & String de 4 dígitos com a senha do sistema\\\hline
        mensagem & String com a última mensagem do sistema\\\hline
        estadoCampainha & Boolean que indica se a campainha está ativa\\\hline
        tentativas & Contador de tentativas de senha falhadas\\\hline
        MAX\_TENTATIVAS & Constante com valor 3\\\hline
    \end{tabular}
    \caption{Atributos da classe ComandoAlarme}
\end{figure}

\textbf{Métodos principais:}

\begin{itemize}
    \item \texttt{armar(senha)} - Arma o sistema se a senha for correta
    \item \texttt{desarmar(senha)} - Desarma o sistema, retorna true se OK
    \item \texttt{mudarSenha(senhaAntiga, senhaNova)} - Altera a senha
    \item \texttt{ativarCampainha()} - Ativa o alarme sonoro
    \item \texttt{processarSensor(tipo, divisao)} - Processa notificação de sensor
\end{itemize}

\subsubsection{Classe Sensor}

\begin{tabular}{|l|p{10cm}|}\hline
    \textbf{Atributos} & \textbf{Descrição}\\\hline
    id & Identificador único do sensor\\\hline
    divisao & Nome da divisão onde o sensor está instalado\\\hline
    estado & Boolean: true=ativo, false=inativo\\\hline
    tipo & Enum TipoSensor (INTRUSAO, MOVIMENTO, FOGO, GAS, INUNDACAO)\\\hline
    comando & Ponteiro para o ComandoAlarme associado\\\hline
\end{tabular}

\textbf{Tipos de Sensores disponíveis:}

\begin{enumerate}
    \item INTRUSAO - Sensores de abertura de portas/janelas
    \item MOVIMENTO - Sensores de movimento interior
    \item FOGO - Sensores de deteção de incêndio
    \item GAS - Sensores de fuga de gás
    \item INUNDACAO - Sensores de inundação
\end{enumerate}

\subsubsection{Classe Simulacao}

Responsável por gerir a fila de eventos e coordenar a simulação. Utiliza uma
fila FIFO (\textit{First-In-First-Out}) para processar eventos sequencialmente.

\begin{tabular}{|l|p{10cm}|}\hline
    \textbf{Atributos} & \textbf{Descrição}\\\hline
    eventos & Queue de objetos Evento a processar\\\hline
    sensores & Vector de ponteiros para sensores\\\hline
    comandoAlarme & Ponteiro para o ComandoAlarme principal\\\hline
\end{tabular}

\subsection{Relações entre Classes}

\begin{figure}[h!]
    \centering
    \textbf{Relacionamentos:}\par\medskip
    \begin{tabular}{lcl}
        Simulacao & $\rightarrow$ & ComandoAlarme (1..1)\\
        Simulacao & $\rightarrow$ & Sensor (1..n)\\
        Sensor & $\rightarrow$ & ComandoAlarme (n..1)
    \end{tabular}
    \caption{Relações entre classes}
\end{figure}

\newpage

\section{Instalação e Compilação}

\subsection{Requisitos}

\begin{itemize}
    \item Compilador C++ com suporte a C++11 ou superior (g++ recomendado)
    \item Sistema operativo: Linux, macOS ou Windows (com MinGW/MSYS2)
    \item Make (para utilizar o Makefile fornecido)
\end{itemize}

\subsection{Estrutura de Ficheiros}

\begin{verbatim}
proj2POO_AnaFreire/
├── ComandoAlarme.hpp     # Declaração da classe
├── ComandoAlarme.cpp     # Implementação
├── Sensor.hpp            # Declaração da classe
├── Sensor.cpp            # Implementação
├── Simulacao.hpp         # Declaração da classe
├── Simulacao.cpp         # Implementação
├── main.cpp              # Programa principal (ficheiro de eventos)
├── mainMenu.cpp          # Programa com menu interativo
├── Makefile              # Ficheiro de compilação
├── eventos.dados         # Ficheiro de eventos de exemplo
├── uml_texto.txt         # Diagrama UML em texto
└── uml_plantuml.txt      # Diagrama UML em PlantUML
\end{verbatim}

\subsection{Compilação}

Para compilar o projeto, utilize os seguintes comandos no terminal:

\subsubsection{Compilar a versão de ficheiro de eventos:}

\begin{lstlisting}[language=bash]
make
\end{lstlisting}

Isto gera o executável \texttt{alarme}.

\subsubsection{Compilar a versão com menu interativo:}

\begin{lstlisting}[language=bash]
make menu
\end{lstlisting}

Isto gera o executável \texttt{alarmeMenu}.

\subsection{Execução}

\begin{lstlisting}[language=bash]
# Executar versao de ficheiro de eventos
./alarme

# Executar versao com menu
./alarmeMenu

# Limpar ficheiros compilados
make clean
\end{lstlisting}

\newpage

\section{Manual de Utilização}

\subsection{Modo Menu Interativo}

Ao executar \texttt{alarmeMenu}, é apresentado um menu com as seguintes opções:

\begin{figure}[h!]
    \centering
    \begin{verbatim}
    ========================================
            MENU SISTEMA DE ALARME
    ========================================
     1. Armar sistema
     2. Desarmar sistema
     3. Ativar sensor
     4. Desativar sensor
     5. Ver estado do sistema
     6. Listar sensores
     7. Mudar senha
     8. Carregar eventos de ficheiro
     9. Desligar campainha
     0. Sair
    ========================================
     Escolha uma opcao:
    \end{verbatim}
    \caption{Menu principal do sistema}
\end{figure}

\subsubsection{Opção 1 - Armar Sistema}

Para armar o sistema, é necessária a senha correta. A senha por omissão é
\textbf{0000}.

\begin{verbatim}
> 1
Introduza a senha: 0000
[ALARME] Sistema armado com sucesso
\end{verbatim}

\subsubsection{Opção 2 - Desarmar Sistema}

Para desarmar, introduza a senha correta. Se errar 3 vezes, a campainha
ativa automaticamente.

\begin{verbatim}
> 2
Introduza a senha: 0000
[ALARME] Sistema desarmado
\end{verbatim}

Erros:

\begin{verbatim}
> 2
Introduza a senha: 1234
[ALARME] Senha incorreta. Tentativas: 1
\end{verbatim}

Se exceder 3 tentativas:

\begin{verbatim}
[ALARME] CAMPANHA ATIVADA - ALARME SONORO
\end{verbatim}

\subsubsection{Opção 3 - Ativar Sensor}

Mostra a lista de sensores e permite ativar um. Quando um sensor é ativado
com o sistema armado, o alarme processa o evento.

\begin{verbatim}
> 3

--- SENSORES ---
Sensor 1: Sala | Movimento | INATIVO
Sensor 2: Cozinha | Gas | INATIVO
Sensor 3: Quarto | Intrusao | INATIVO
Sensor 4: Casa de Banho | Inundacao | INATIVO
Sensor 5: Entrada | Intrusao | INATIVO
ID do sensor a ativar: 1
[SENSOR 1] Ativado na divisao Sala
\end{verbatim}

Se o sistema estiver armado, é mostrada uma mensagem de alerta.

\subsubsection{Opção 4 - Desativar Sensor}

Desativa um sensor específico.

\begin{verbatim}
> 4
ID do sensor a desativar: 1
[SENSOR 1] Desativado
\end{verbatim}

\subsubsection{Opção 5 - Ver Estado do Sistema}

Mostra o estado atual do alarme.

\begin{verbatim}
> 5

--- ESTADO DO SISTEMA ---
Sistema: ARMADO
Campainha: DESATIVA
Tentativas: 0/3
Mensagem: Sistema armado com sucesso
\end{verbatim}

\subsubsection{Opção 6 - Listar Sensores}

Lista todos os sensores com as suas características.

\begin{verbatim}
> 6

--- SENSORES ---
Sensor 1: Sala | Movimento | ATIVO
Sensor 2: Cozinha | Gas | INATIVO
Sensor 3: Quarto | Intrusao | INATIVO
Sensor 4: Casa de Banho | Inundacao | INATIVO
Sensor 5: Entrada | Intrusao | INATIVO
\end{verbatim}

\subsubsection{Opção 7 - Mudar Senha}

Altera a senha do sistema. A nova senha deve ter exatamente 4 dígitos.

\begin{verbatim}
> 7
Senha atual: 0000
Nova senha (4 digitos): 1234
[ALARME] Senha alterada com sucesso
\end{verbatim}

Erro se senha atual estiver errada:

\begin{verbatim}
[ALARME] Senha atual incorreta
\end{verbatim}

Erro se nova senha não tiver 4 dígitos:

\begin{verbatim}
[ALARME] Nova senha deve ter 4 digitos
\end{verbatim}

\subsubsection{Opção 8 - Carregar Eventos de Ficheiro}

Carrega e executa eventos de um ficheiro de texto.

\begin{verbatim}
> 8
Nome do ficheiro de eventos: eventos.dados
[SIMULACAO] A iniciar simulacao...
[SIMULACAO] Carregados 10 eventos
[SIMULACAO] Eventos carregados com sucesso

========================================
       INICIO DA SIMULACAO
========================================
...
\end{verbatim}

\subsubsection{Opção 9 - Desligar Campainha}

Desativa manualmente a campainha.

\begin{verbatim}
> 9
[ALARME] Campainha desativada
\end{verbatim}

\subsection{Modo Ficheiro de Eventos}

O programa \texttt{alarme} executa automaticamente os eventos definidos no
ficheiro \texttt{eventos.dados}.

\newpage

\section{Formato do Ficheiro de Eventos}

\subsection{Estrutura do Ficheiro}

O ficheiro de eventos é um ficheiro de texto onde cada linha representa um
evento a ser processado. Os eventos são executados por ordem de chegada
(FIFO - First-In-First-Out).

\subsection{Comandos Disponíveis}

\begin{table}[h!]
    \centering
    \begin{tabular}{|l|l|p{7cm}|}\hline
        \textbf{Comando} & \textbf{Parâmetros} & \textbf{Descrição}\\\hline
        ARMAR & <senha> & Arma o sistema com a senha indicada\\\hline
        DESARMAR & <senha> & Desarma o sistema\\\hline
        MUDARSENHA & <senha\_antiga> <senha\_nova> & Altera a senha\\\hline
        ATIVAR & <id\_sensor> & Ativa o sensor especificado\\\hline
        DESATIVAR & <id\_sensor> & Desativa o sensor especificado\\\hline
    \end{tabular}
    \caption{Comandos disponíveis no ficheiro de eventos}
\end{table}

\subsection{Exemplo de Ficheiro}

\begin{lstlisting}[caption={ficheiro eventos.dados}]
ARMAR 0000
ATIVAR 1
DESARMAR 0000
MUDARSENHA 0000 1234
ATIVAR 2
ATIVAR 3
DESARMAR 1234
ATIVAR 5
MUDARSENHA 1234 9999
DESARMAR 9999
\end{lstlisting}

\subsection{Resultado da Execução}

\begin{verbatim}
========================================
       INICIO DA SIMULACAO
========================================

--- Processando evento: ARMAR 0000 ---
[ALARME] Sistema armado com sucesso

--- Processando evento: ATIVAR 1 ---
[SENSOR 1] Ativado na divisao Sala

--- Processando evento: DESARMAR 0000 ---
[ALARME] Sistema desarmado

--- Processando evento: MUDARSENHA 0000 1234 ---
[ALARME] Senha alterada com sucesso
...
\end{verbatim}

\subsection{Exemplo: Teste de Tentativas Erradas}

\begin{lstlisting}[caption={ficheiro para testar campainha}]
ARMAR 0000
ATIVAR 1
DESARMAR 1111
DESARMAR 2222
DESARMAR 3333
DESARMAR 0000
\end{lstlisting}

Resultado:

\begin{verbatim}
--- Processando evento: DESARMAR 1111 ---
[ALARME] Senha incorreta. Tentativas: 1

--- Processando evento: DESARMAR 2222 ---
[ALARME] Senha incorreta. Tentativas: 2

--- Processando evento: DESARMAR 3333 ---
[ALARME] CAMPANHA ATIVADA - ALARME SONORO
[ALARME] Senha incorreta. Tentativas: 3

--- Processando evento: DESARMAR 0000 ---
[ALARME] Sistema desarmado
\end{verbatim}

\newpage

\section{Diagramas UML}

\subsection{Diagrama de Classes (Texto)}

\begin{verbatim}
CLASSE ComandoAlarme
--------------------
+---------------+
| ComandoAlarme  |
+---------------+
| - estado       |
| - senha        |
| - mensagem     |
| - campainha    |
| - tentativas   |
+---------------+
| + armar()      |
| + desarmar()   |
| + mudarSenha() |
| + ativaCamp()  |
| + getEstado() |
+---------------+

CLASSE Sensor
-------------
+---------------+
| Sensor        |
+---------------+
| - id          |
| - divisao     |
| - estado      |
| - tipo        |
| - comando*    |
+---------------+
| + ativar()    |
| + desativar() |
| + getId()     |
+---------------+

CLASSE Simulacao
-----------------
+--------------------+
| Simulacao          |
+--------------------+
| - eventos: queue   |
| - sensores: vector |
| - comandoAlarme*   |
+--------------------+
| + iniciar()        |
| + executar()       |
| + obtemSensor()    |
+--------------------+
\end{verbatim}

\subsection{Diagrama PlantUML}

O ficheiro \texttt{uml\_plantuml.txt} contém o código PlantUML que pode ser
convertido para imagem utilizando a ferramenta PlantUML ou diversos sites online.

\begin{verbatim}
@startuml
class ComandoAlarme {
    - estado : bool
    - senha : string
    - mensagem : string
    - estadoCampainha : bool
    - tentativas : int
    + armar(senha : string)
    + desarmar(senha : string) : bool
    + mudarSenha(senhaAntiga, senhaNova : string) : bool
    + ativarCampainha()
}

class Sensor {
    - id : int
    - divisao : string
    - estado : bool
    - tipo : TipoSensor
    - comando : ComandoAlarme*
    + ativar()
    + desativar()
    + getId() : int
}

class Simulacao {
    - eventos : queue<Evento>
    - sensores : vector<Sensor*>
    - comandoAlarme : ComandoAlarme*
    + iniciar(ficheiro : string) : bool
    + executar()
    + obtemSensor(id : int) : Sensor*
}

Simulacao "1" o-- "1" ComandoAlarme
Simulacao "1" o-- "1..n" Sensor
Sensor "n" --> "1" ComandoAlarme
@enduml
\end{verbatim}

\newpage

\section{Testes e Validação}

\subsection{Cenários de Teste}

\subsubsection[Teste 1: Armar e desarmar normal]
\textbf{Teste 1: Armar e desarmar normal}

\begin{enumerate}
    \item Armar com senha correta (0000)
    \item Verificar que sistema fica armado
    \item Desarmar com senha correta
    \item Verificar que sistema fica desarmado
\end{enumerate}

\textbf{Resultado esperado:} Sistema responde corretamente a ambos os comandos.

\subsubsection[Teste 2: Tentativas erradas]
\textbf{Teste 2: Tentativas erradas}

\begin{enumerate}
    \item Armar sistema
    \item Tentar desarmar com senha errada 3 vezes
    \item Verificar que campainha ativa após 3ª tentativa
    \item Desarmar com senha correta
    \item Verificar que campainha desativa
\end{enumerate}

\textbf{Resultado esperado:} Campainha ativa após 3 tentativas falhadas.

\subsubsection[Teste 3: Alteração de senha]
\textbf{Teste 3: Alteração de senha}

\begin{enumerate}
    \item Alterar senha de 0000 para 1234
    \item Desarmar com nova senha (1234)
    \item Verificar que senha antiga (0000) já não funciona
\end{enumerate}

\textbf{Resultado esperado:} Senha alterada com sucesso e validação funciona.

\subsection{Testes de Sensores}

\begin{enumerate}
    \item Criar sensores em diferentes divisões
    \item Ativar cada tipo de sensor
    \item Verificar que notificação é enviada ao ComandoAlarme
    \item Confirmar que sensores podem ser desativados
\end{enumerate}

\newpage

\section{Conclusão}

Este projeto de Sistema de Alarme Doméstico demonstra a aplicação de conceitos
fundamentais de Programação Orientada aos Objetos em C++, incluindo:

\subsection{Conceitos POO Implementados}

\begin{itemize}
    \item \textbf{Encapsulamento}: Atributos privados com métodos públicos de acesso
    \item \textbf{Composição}: Classe Simulacao contém objetos ComandoAlarme e Sensor
    \item \textbf{Associação}: Sensores conhece o ComandoAlarme ao qual estão ligados
    \item \textbf{Gestão de memória}: Uso correto de ponteiros
    \item \textbf{STL}: Utilização de queue, vector e string
\end{itemize}

\subsection{Limitações e Trabalho Futuro}

\begin{itemize}
    \item Adicionar herança com classe base Componente (como na correção do professor)
    \item Implementar mais tipos de eventos
    \item Adicionar timestamps aos eventos
    \item Persistência de dados em ficheiros
    \item Interface gráfica (GUI)
    \item Suporte a múltiplos utilizadores
\end{itemize}

\subsection{Notas de Desenvolvimento}

O sistema foi desenvolvido e testado utilizando:

\begin{itemize}
    \item Compilador: g++ com C++11
    \item Sistema operativo: Linux
    \item Controlo de versões: Git
    \item Repositório remoto: GitHub
\end{itemize}

\newpage

\section{Referências}

\begin{thebibliography}{9}

    \bibitem{c++standard}
    ISO/IEC JTC1/SC22/WG21 - The C++ Programming Language.
    \textit{C++11, C++14, C++17, C++20 Standards}.

    \bibitem{stl}
    Standard Template Library Documentation.
    \textit{cppreference.com}.

    \bibitem{uml}
    UML (Unified Modeling Language) Documentation.
    \textit{OMG UML Specification}.

    \bibitem{plantuml}
    PlantUML - Open Source UML Diagram Generator.
    \textit{www.plantuml.com}.

\end{thebibliography}

\newpage

\section{Anexos}

\subsection{Lista de Ficheiros do Projeto}

\begin{table}[h!]
    \centering
    \begin{tabular}{|l|p{8cm}|}\hline
        \textbf{Ficheiro} & \textbf{Descrição}\\\hline
        ComandoAlarme.hpp & Declaração da classe ComandoAlarme\\\hline
        ComandoAlarme.cpp & Implementação da classe ComandoAlarme\\\hline
        Sensor.hpp & Declaração da classe Sensor\\\hline
        Sensor.cpp & Implementação da classe Sensor\\\hline
        Simulacao.hpp & Declaração da classe Simulacao\\\hline
        Simulacao.cpp & Implementação da classe Simulacao\\\hline
        main.cpp & Programa principal (execução por ficheiro)\\\hline
        mainMenu.cpp & Programa com menu interativo\\\hline
        Makefile & Ficheiro de compilação\\\hline
        eventos.dados & Ficheiro de eventos de exemplo\\\hline
        uml\_texto.txt & Diagrama UML em formato texto\\\hline
        uml\_plantuml.txt & Diagrama UML em PlantUML\\\hline
    \end{tabular}
    \caption{Ficheiros do projeto}
\end{table}

\subsection{Comandos Make}

\begin{table}[h!]
    \centering
    \begin{tabular}{|l|p{8cm}|}\hline
        \textbf{Comando} & \textbf{Descrição}\\\hline
        make & Compila a versão de ficheiro de eventos\\\hline
        make menu & Compila a versão com menu interativo\\\hline
        make run & Compila e executa a versão de ficheiro\\\hline
        make clean & Remove os executáveis compilados\\\hline
    \end{tabular}
    \caption{Comandos do Makefile}
\end{table}

\end{document}